% =============================================================
%  Chapter 2 — Related Work (shortened draft)
% =============================================================
%    §2.1  LLM code generation and backend evaluation benchmarks
%    §2.2  Self-refinement, agentic coding, and verification-gated
%           generation
%    §2.3  Deployment configuration and learned cluster management
%           (incl. LLM-generated infrastructure as code)
%    §2.4  LLM-driven performance optimization
%    §2.5  Positioning of this work
% =============================================================

\chapter{Related Work}
\label{ch:related-work}


% ----- §2.1  Benchmarks -----
\section{LLM Code Generation and Backend Evaluation Benchmarks}
\label{sec:rw-benchmarks}

Evaluation of large language models for software engineering has moved from
single-function synthesis~\cite{chen2021humaneval,austin2021mbpp} to
repository-level issue resolution~\cite{jimenez2024swebench} to complete,
security-checked backend applications~\cite{vero2025baxbench}
(Section~\ref{sec:background-codegen}).  CWEval~\cite{peng2025cweval} and
SecRepoBench~\cite{shen2026secrepobench} extend the security dimension with
outcome-driven test oracles and repository-scale, fuzz-tested code-agent
evaluation, respectively.  CodeFlowBench~\cite{wang2025codeflowbench} extends
the correctness dimension to \emph{multi-turn} construction, where later
functionality reuses earlier turns' output, and finds significant
degradation as dependency complexity grows.

Together, these benchmarks characterise whether generated software is
correct, secure, or constructible over multiple turns.  All of them execute
programs in isolated containers, without orchestrated deployment, shared
cluster resources, or load testing (Section~\ref{sec:background-containerization}).
None measures \emph{sustainable throughput} when an application is deployed
as a replicated service on Kubernetes, which is the gap our work addresses.


% ----- §2.2  Self-refinement, agentic coding, and verification -----
\section{Self-Refinement, Agentic Coding, and Verification-Gated Generation}
\label{sec:rw-agentic}

Section~\ref{sec:background-iterative-refinement} introduces the general
propose-evaluate-revise pattern our own method instantiates.  Two lines of
prior work instantiate it for code, with different feedback sources.
Reflexion~\cite{shinn2023reflexion} and Self-Refine~\cite{madaan2023selfrefine}
use an LLM's own linguistic self-critique, an agent verbalises what went
wrong, or critiques its own output, and carries that text into its next
attempt, with no external verifier.  Specialised for software engineering,
SWE-agent~\cite{yang2024sweagent} and OpenHands~\cite{wang2024openhands}
instead close the loop with compiler and test signals from a sandboxed
shell, the latter now underlying much of the field's pluggable multi-agent
tooling.

A related but distinct line constrains what an LLM may output, rather than
iterating on what it has already produced.
Synchromesh~\cite{poesia2022synchromesh} restricts token sampling to
grammar- and type-correct programs; Clover~\cite{sun2023clover} and
VERGE~\cite{singh2026verge} instead check code, specifications, or
decomposed logical claims for mutual consistency using formal tools or SMT
solvers.  All three target small, well-specified problems and have not been
shown to scale to BaxBench-sized applications.

Our agent instead revises one of two structured artefacts, application code
or a constrained deployment specification, using quantitative signals from
an actual multi-node Kubernetes deployment under load: goodput, latency
percentiles, and resource utilisation.  No self-critique step or formal
proof is required, since the load test itself supplies an objective, numeric
verifier that neither line above uses, and our correctness gate, BaxBench
functional tests plus schema-validated specifications, is coarser but more
scalable than constrained decoding or SMT verification.


% ----- §2.3  Cluster tuning -----
\section{Deployment Configuration and Learned Cluster Management}
\label{sec:rw-cluster}

Self-management, a system observing its own state and adjusting itself
toward an operator's goals, predates LLM agents.  Kephart and Chess's
autonomic computing~\cite{kephart2003autonomic} frames this as a
monitor-analyse-plan-execute loop; our own iteration loop (bench, feedback,
decide, refine) is an instance of it, with an LLM standing in for the fixed
control law.  In Kubernetes~\cite{burns2016borg,kubernetes}
(Sections~\ref{sec:background-kubernetes}--\ref{sec:background-deployment-levers}),
day-to-day tuning is typically \emph{reactive}: the Horizontal and Vertical
Pod Autoscalers~\cite{kubernetes-hpa} adjust replica count or per-pod
resources from live metrics within a fixed template, without reasoning
about database topology, connection pooling, or pod placement.

Several systems instead learn resource-management policy from data while
holding the application and its manifest fixed.
Autopilot~\cite{rzadca2020autopilot} recommends horizontal and vertical
scaling from historical utilisation; FIRM~\cite{qiu2020firm} localises and
automatically reprovisions the resource causing a service-level-objective
violation; Decima~\cite{mao2019decima} learns cluster scheduling policy with
deep reinforcement learning; OtterTune~\cite{vanaken2017ottertune} applies
the same idea to DBMS configuration knobs directly analogous to our
connection-pooling and database-topology fields.  None of the four revises
application code or a deployment manifest for arbitrary HTTP services, and
none is driven by an LLM.

Glia~\cite{hamadanian2025glia} is the closest prior work in combining LLM
reasoning with measured performance feedback for infrastructure behaviour: a
multi-agent architecture that autonomously designs request-routing,
scheduling, and autoscaling mechanisms for a GPU cluster serving LLM
inference.  It targets inference-serving mechanism design, generating new
algorithms, rather than filling a constrained deployment specification for
general backends.

Closer to the \emph{specification} side of our framework,
\textsc{IaCGen}~\cite{zhang2026iacgen} and TerraFormer~\cite{jana2026terraformer}
apply the same iterative-agent pattern directly to infrastructure code
(surveyed by Srivatsa~et~al.~\cite{srivatsa2024iacsurvey}): a
fail-learn-refine-succeed loop against a live deployment simulation
(\textsc{IaCGen}, raising first-attempt success from 20--30\% to over 90\%
within 25~attempts), or against policy-guided verifier feedback
(TerraFormer).  Both stop at \emph{deployability}, a binary, did-it-provision
signal, rather than the throughput or latency the provisioned system
sustains, and neither touches the application code being deployed.  Our
deploy probe plays exactly this role, but as a precondition for our loop
rather than its objective: only after a candidate is schedulable and ready
do we measure and refine for goodput.

No prior system combines both limits: learned resource managers (FIRM,
Autopilot, Decima, OtterTune) hold application and specification fixed and
optimise policy instead; IaC-generation agents (\textsc{IaCGen},
TerraFormer) refine a specification but stop at deployability.  Refining a
constrained deployment specification \emph{and} application code, driven by
runtime load-test goodput rather than a pass/fail deployability check, is
the problem our framework formalises.


% ----- §2.4  Performance optimization -----
\section{LLM-Driven Performance Optimization}
\label{sec:rw-llm-perf}

Automatically searching for a faster configuration of a fixed program
predates LLMs: OpenTuner~\cite{ansel2014opentuner} runs an ensemble of
search techniques over a program's parameter space, reallocating budget
toward whichever technique is currently improving a measured objective
fastest.  It shares our premise that a measured objective should drive the
next candidate, but searches a numeric parameter space rather than source
code or a structured deployment specification.

More recent work asks whether LLMs specifically can improve
\emph{performance}, not merely correctness, given execution feedback.
Shypula~et~al.~\cite{shypula2024pie} curate PIE, over 77{,}000~human
performance-improving edit pairs from competitive programming, and show that
LLMs given feedback from a timing harness can learn to produce measurably
faster code, including algorithmic rewrites.  SWE-Perf~\cite{he2025sweperf}
moves the same question to repository scale, 140~instances derived from real
performance-improving pull requests, and finds a substantial gap between
current LLMs and expert optimisers.  Both are largely \emph{one-shot}:
neither incorporates iterative runtime feedback from a system executing
under load.

Agentic systems have begun to close that loop with measurement.
VibeServe~\cite{kamahori2025vibeserve} synthesises bespoke LLM serving
runtimes through a two-level multi-agent loop, an outer planner tracking
optimisation history, an inner loop separating implementation, correctness
judging, and performance evaluation, matching vLLM on standard Llama serving
and achieving large speedups on non-standard hardware-workload combinations.
Liu~et~al.~\cite{liu2025jitsystems}'s \emph{Just-in-Time Systems}
(\textsc{Jitskit}) iterate an LLM agent through planning, coding, and
critic-guided reflection to synthesise entire key-value stores that
outperform FASTER, RocksDB, and Redis on all 18~configurations tested.

Our framework shares this iterative structure, propose, gate on correctness,
benchmark, feed diagnostics forward, but at a different level of the stack:
PIE and SWE-Perf optimise \emph{source code} in isolation; VibeServe and
Jitskit synthesise \emph{system implementations}; we optimise
\emph{Kubernetes deployment configuration} (and optionally application code)
for BaxBench backends, using \emph{goodput} as the objective under an
adaptive load profile.


% ----- §2.5  Positioning -----
\section{Positioning of This Work}
\label{sec:rw-positioning}

Table~\ref{tab:rw-comparison} summarises how representative systems relate to
our framework across evaluation target, feedback structure, and deployment
environment.  Benchmarks through CodeFlowBench establish correctness,
security, and multi-turn code quality but not cluster performance.
General-purpose and code-specialised refinement loops, and verification-based
reliability methods, establish that LLMs can improve or constrain their own
output, but none is driven by a quantitative infrastructure metric.
Classical and learned cluster or database tooling, and LLM-driven IaC
generation, address infrastructure tuning or deployability without
benchmarking full LLM-generated backends under adaptive load.  Glia is
closest in spirit, combining LLM reasoning with measured feedback, but for
inference-serving mechanism design.  SWE-Perf, PIE, VibeServe, and Jitskit
demonstrate LLM-driven performance optimisation at the code or
implementation level.

Our thesis occupies the intersection these lines leave open: \emph{iterative
deployment optimisation} of complete, LLM-generated backend applications on a
real Kubernetes cluster, with the agent controlling both application code and
a validated deployment specification, scored each iteration by measured
goodput rather than static tests, self-critique, a deployability check, or a
fixed benchmark application alone.

\begin{table}[t]
  \centering
  \caption{Comparison of related approaches (selected dimensions).}
  \label{tab:rw-comparison}
  \small
  \begin{tabular}{@{}l c c c c c c@{}}
    \toprule
    & \textbf{BaxBench} & \textbf{Glia} & \textbf{IaCGen} & \textbf{SWE-Perf} &
      \textbf{VibeServe} & \textbf{This work} \\
    \midrule
    Primary focus
      & Func.+sec.
      & Mechanism design
      & Deployability
      & Code perf.
      & Serving runtime
      & Deploy perf. \\
    Iterative loop
      & No
      & Yes
      & Yes
      & No
      & Yes
      & Yes \\
    Runtime load test
      & No
      & Yes
      & No
      & Perf.\ tests
      & Yes
      & Yes \\
    Environment
      & Docker
      & GPU cluster
      & Multi-cloud
      & Local/CI
      & Single GPU
      & Multi-node K8s \\
    Agent controls
      & Code
      & Algorithms
      & IaC spec
      & Code patches
      & System code
      & Code + spec \\
    \bottomrule
  \end{tabular}
\end{table}
